% =============================================================================
% Module 1a: Solutions to Exercises
% =============================================================================
% This file is conditionally included based on \ifsolutions
% =============================================================================

\section{Solutions to Exercises}
\label{sec:sr_solutions}

\noindent
Full worked solutions are also available in the Mathematica script
\solnlink{01a_Special_Relativity/problems\_01a.wl}{problems\_01a.wl},
which verifies each answer symbolically.

% ----- Exercise 1.1 -----
\subsection*{Exercise~\ref{ex:interval_classification}: Interval Classification}

We compute $(\Delta s)^2 = -(\Delta t)^2 + (\Delta x)^2 + (\Delta y)^2 + (\Delta z)^2$:
\begin{enumerate}[label=(\alph*)]
    \item $(\Delta s)^2 = -(2)^2 + (1)^2 + (1)^2 + (0)^2 = -4 + 1 + 1 = -2 < 0$
          $\Rightarrow$ \textbf{Timelike}.  A massive particle can travel between these events.

    \item $(\Delta s)^2 = -(1)^2 + (1)^2 + (0)^2 + (0)^2 = 0$
          $\Rightarrow$ \textbf{Null}.  A light ray connects these events.

    \item $(\Delta s)^2 = -(3)^2 + (1)^2 + (2)^2 + (2)^2 = -9 + 1 + 4 + 4 = 0$
          $\Rightarrow$ \textbf{Null}.  Again lightlike: $|\Delta\mathbf{x}|^2 = 9 = (\Delta t)^2$.
\end{enumerate}

% ----- Exercise 1.2 -----
\subsection*{Exercise~\ref{ex:boost_composition}: Boost Composition}

The boost matrix for rapidity $\phi$ is
$\Lambda(\phi) = \begin{psmallmatrix} \cosh\phi & -\sinh\phi \\ -\sinh\phi & \cosh\phi \end{psmallmatrix}$
(showing only the $t$--$x$ block).  Computing the product:
\begin{align*}
    \Lambda(\phi_1)\,\Lambda(\phi_2)
    &= \begin{pmatrix}
        \cosh\phi_1\cosh\phi_2 + \sinh\phi_1\sinh\phi_2 &
        -\cosh\phi_1\sinh\phi_2 - \sinh\phi_1\cosh\phi_2 \\
        -\sinh\phi_1\cosh\phi_2 - \cosh\phi_1\sinh\phi_2 &
        \sinh\phi_1\sinh\phi_2 + \cosh\phi_1\cosh\phi_2
    \end{pmatrix} \\
    &= \begin{pmatrix}
        \cosh(\phi_1+\phi_2) & -\sinh(\phi_1+\phi_2) \\
        -\sinh(\phi_1+\phi_2) & \cosh(\phi_1+\phi_2)
    \end{pmatrix}
    = \Lambda(\phi_1 + \phi_2) \,,
\end{align*}
using the hyperbolic addition formulas.

For the velocity addition formula, we use $v = \tanh\phi$:
\begin{equation*}
    v = \tanh(\phi_1 + \phi_2)
    = \frac{\tanh\phi_1 + \tanh\phi_2}{1 + \tanh\phi_1\tanh\phi_2}
    = \frac{v_1 + v_2}{1 + v_1 v_2} \,.
\end{equation*}
Note: for $v_1 = v_2 = 1$ (speed of light), $v = 2/2 = 1$, confirming
that the speed of light is invariant.

% ----- Exercise 1.3 -----
\subsection*{Exercise~\ref{ex:four_velocity_norm}: Four-Velocity Norm}

The four-velocity is $U^\mu = \gamma(1, v^1, v^2, v^3)$ where
$\gamma = (1 - v^2)^{-1/2}$ and $v^2 = (v^1)^2 + (v^2)^2 + (v^3)^2$.
Then:
\begin{align*}
    U^\mu U_\mu &= \eta_{\mu\nu} U^\mu U^\nu
    = \gamma^2\bigl[-(1)^2 + (v^1)^2 + (v^2)^2 + (v^3)^2\bigr] \\
    &= \gamma^2(-1 + v^2)
    = \frac{-(1 - v^2)}{1 - v^2}
    = \boxed{-1} \,.
\end{align*}

% ----- Exercise 1.4 -----
\subsection*{Exercise~\ref{ex:photon_energy}: Photon Energy (Relativistic Doppler)}

\textbf{(a)} A photon with energy $E$ propagating at angle $\theta$ to the
$x$-axis has four-momentum $p^\mu = E(1, \cos\theta, \sin\theta, 0)$.
Applying the boost $\Lambda$ along $x$:
\begin{equation*}
    E' = p^{0'} = \gamma(p^0 - v\,p^1) = \gamma E(1 - v\cos\theta) \,.
\end{equation*}

\textbf{(b)} Special cases:
\begin{itemize}
    \item $\theta = 0$ (photon along $+x$, same direction as boost):
    $E' = \gamma E(1-v) = E\sqrt{\frac{1-v}{1+v}}$.
    This is a \textbf{redshift} ($E' < E$).  The observer recedes from the source.

    \item $\theta = \pi/2$ (transverse):
    $E' = \gamma E = E/\sqrt{1-v^2}$.
    This is the \textbf{transverse Doppler effect} --- a purely relativistic
    blueshift with no classical analogue, arising from time dilation.
\end{itemize}

% ----- Exercise 1.5 -----
\subsection*{Exercise~\ref{ex:em_tensor}: Electromagnetic Field Strength Tensor}

\textbf{(a)} By inspection, $F_{\mu\nu} = -F_{\nu\mu}$ (the matrix is
antisymmetric: the diagonal is zero, and off-diagonal elements differ by a sign).

\textbf{(b)} Raising both indices with the metric:
$F^{\mu\nu} = \eta^{\mu\alpha}\eta^{\nu\beta}F_{\alpha\beta}$.
In matrix form: $F^{\mu\nu} = \eta \cdot F_{\mu\nu} \cdot \eta$.
The result is:
\[
    F^{\mu\nu} =
    \begin{pmatrix}
    0 & E_1 & E_2 & E_3 \\
    -E_1 & 0 & B_3 & -B_2 \\
    -E_2 & -B_3 & 0 & B_1 \\
    -E_3 & B_2 & -B_1 & 0
    \end{pmatrix}
\]
The electric field components flip sign (from raising the time index),
while the magnetic components are unchanged.

\textbf{(c)} The Lorentz scalar:
\begin{align*}
    F_{\mu\nu}F^{\mu\nu}
    &= 2\bigl(-E_1^2 - E_2^2 - E_3^2 + B_1^2 + B_2^2 + B_3^2\bigr) \\
    &= 2(|\mathbf{B}|^2 - |\mathbf{E}|^2) \,.
\end{align*}
This is a Lorentz invariant: all inertial observers agree on the value of
$|\mathbf{B}|^2 - |\mathbf{E}|^2$.
